\par\bigskip
\begin{center}\pdfbookmark[2]{7. Secular and periodic perturbations}{ceres-007}[7.]\end{center}
\noindent\hspace*{\parindent}From the formulas above we now obtain

I) Secular increase in eccentricity \(=\tfrac12e'a\mathfrak C^{(0)}nt\sin(\tilde\omega'-\tilde\omega)\). [Now] \(\mathfrak C^{(0)}=\tfrac12Q''\), so the increase \(=\tfrac14e'aQ''nt\sin(\tilde\omega'-\tilde\omega)\)
\[
\begin{array}{cc}
\text{For Ceres:}&\text{[For Jupiter:]}\\[.5ex]
\tilde\omega=352^{0}\,58'\,0''&\tilde\omega'=191^{0}\,10'\,20''\\[.5ex]
\log n\ \text{[daily]}=2{,}886\,2959&\log e'=8{,}681\,3317\,[-10]
\end{array}
\]
[and from this we find
\[
\text{annual decrease in eccentricity]}\quad0{,}00000\,59089.
\]

II) Secular increase in the longitude of aphelion:
\[
=\frac a4\left(B^{(0)}-2\mathfrak B^{(0)}-2\frac{e'}e\mathfrak C^{(0)}\cos(\tilde\omega'-\tilde\omega)\right)nt.
\]
Now
\[
B^{(0)}-2\mathfrak B^{(0)}=Q',
\]
so that increase is
\[
=\frac a4\left(Q'-\frac{e'}eQ''\cos(\tilde\omega'-\tilde\omega)\right)nt.
\]

For Ceres, annual motion of the apsides \(=+70''{,}1503.\)

After these secular changes in eccentricity and aphelion, the equations underlined above in red\refstepcounter{footnote}\hyperlink{ceres7-note}{[*)]}\begingroup
\makeatletter
\long\def\@makefntext#1{\noindent[*) #1}
\makeatother
\footnotetext{\hypertarget{ceres7-note}{}Page 389, where they were also underlined in the printed edition.]}
\endgroup{} drop out entirely.

III) The constant part of \(\delta r\) is \(=[-aa(2g+\mathfrak A^{(0)})=]+\tfrac16aaB^{(0)}\)
\[
\text{[thus]}\quad=-0{,}000\,9475.
\]

IV) The coefficient of \(\cos iD\) [in \(\delta r\)] becomes (because \(i\) previously took both positive and negative values and is now restricted to positive values)
\[
\left[\frac{2aa\mathfrak A^{(i)}}{hh-1}=\right]\frac{2aa}{hh-1}\left(\frac ihA^{(i)}+\tfrac12B^{(i)}\right)\mathpunct{\cdot}
\]

\hypertarget{native-398}{}[In our case, \(\log n'=2{,}475\,8562\),] we obtain \(\log\theta=9{,}589\,5603\) [and with it the part of \(\delta r\):]
\[
\begin{aligned}
&+0{,}001\,0304\cos D&&-0{,}000\,0350\cos5D\\
&-0{,}003\,8023\cos2D&&-0{,}000\,0128\cos6D\\
&-0{,}000\,4206\cos3D&&-0{,}000\,0050\cos7D.\\
&-0{,}000\,1077\cos4D
\end{aligned}
\]

V) In \(\delta v\) the coefficient of \(\sin iD\) becomes
\[
\begin{aligned}
&-\frac{4ai(1-\theta)}{ii(1-\theta)^2-1}\left(\frac1{1-\theta}A^{(i)}+\tfrac12B^{(i)}\right)\\
&+\frac a{i(1-\theta)}\left(\frac{3A^{(i)}}{1-\theta}+2B^{(i)}\right)
\end{aligned}
\]
[and so in our case the corresponding terms are:]
\[
\begin{aligned}
&-231''{,}94\sin D&&+3''{,}05\sin5D\\
&+496{,}68\sin2D&&+1{,}07\sin6D\\
&+44{,}15\sin3D&&+0{,}41\sin7D.\\
&+10{,}07\sin4D
\end{aligned}
\]

VI) The following analysis serves to find the [terms depending on] \(M\) and \(M'+D\). [So far, the constants of integration contained in equation 11 have been left out of account in \(\delta v\). They may be chosen arbitrarily, since they combine with the elliptic elements \(e\) and \(\tilde\omega\), and are to be determined so that the terms depending on \(M\) and \(D+M'=M+\tilde\omega-\bar{\tilde\omega}'\) vanish in \(\delta v\). If the constants of integration \(f_1\) and \(f_1'\) are introduced so that the constant of \(\int rQ\sin V\mathord{.}\,n\,\mathrm{d}t\) equals \(aaef_1+aae'f_1'\cos(\tilde\omega-\tilde\omega')\) and that arising from \(\int rQ\cos V\mathord{.}\,n\,\mathrm{d}t\) equals \(aae'f_1'\sin(\tilde\omega-\tilde\omega')\), then, neglecting the second powers of the eccentricities, the corresponding part of \(\delta r\) is
\[
\alpha)\qquad aaef_1\cos M+aae'f_1'\cos(M+\tilde\omega-\tilde\omega')
\]
and this gives rise in \(\delta v\) to the part
\[
\beta)\qquad-2aef_1\sin M-2ae'f_1'\sin(M+\tilde\omega-\tilde\omega').]
\]

Let the [complete set of \(M\)- and \(D+M'\)-dependent] terms, including those [\(\alpha\)] arising from the constants, be
\hypertarget{native-399}{}
\[
\text{in }\delta r\ .\, .\, .\, .\quad aaef\cos M+aae'f'\cos(M+\tilde\omega-\tilde\omega')
\]
[so that
\[
\begin{aligned}
f&=f_1-\tfrac1{24}B^{(0)}-\tfrac14\mathfrak B^{(0)}\\
f'&=f_1'-\tfrac14\mathfrak C^{(0)}.
\end{aligned}
\]

This makes the entire part depending on the two arguments
\[
\begin{aligned}
\text{in }\delta v\ .\, .\, .\, .\quad&-2a\{f-\tfrac16B^{(0)}-\tfrac12\mathfrak B^{(0)}+\tfrac14\beta^{(0)}\}e\sin M\\
&-2a\{f'-\tfrac12\mathfrak C^{(0)}+\tfrac14\gamma^{(0)}\}e'\sin(M+\tilde\omega-\tilde\omega').]
\end{aligned}
\]
Determine \(f,f'\) so that the sum of these quantities vanishes in \(\delta v\). This leaves
\[
\begin{aligned}
\text{in }\delta r\ .\, .\, .\, .\quad&aae(\tfrac12\mathfrak B^{(0)}-\tfrac14\beta^{(0)}+\tfrac16B^{(0)})\cos M\\
&+aae'(\tfrac12\mathfrak C^{(0)}-\tfrac14\gamma^{(0)})\cos(M+\tilde\omega-\tilde\omega')
\end{aligned}
\]
or
\[
\begin{aligned}
&aae(\tfrac5{12}B^{(0)}+\tfrac14C^{(0)})\cos M\\
&+aae'(\tfrac14A'-\tfrac14C')\cos(M+\tilde\omega-\tilde\omega')
\end{aligned}
\]
[and numerically]
\[
-0{,}000\,0526\cos M\qquad+0{,}000\,0141\cos(M+\tilde\omega-\tilde\omega').
\]

[These two terms can be combined into one. If we temporarily denote them by
\[
b_1\cos M+b_2\cos(M+\tilde\omega-\tilde\omega')
\]
and determine the constants \(c\) and \(C\) from the equations
\[
\begin{aligned}
c\cos C&=b_1+b_2\cos(\tilde\omega-\tilde\omega')\\
c\sin C&=b_2\sin(\tilde\omega-\tilde\omega'),
\end{aligned}
\]
they take the form
\[
c\cos(M+C),
\]
numerically] \(-0{,}000\,0633\cos(M+25^{0}\,56'\,46'').\)

VII) In \(\delta r\) the coefficient of \(\cos(iD+M)\), [which] \(aaeE^{(i)}\) [may denote], is given by
\[
\begin{aligned}
aaeE^{(i)}&=aae\left\{\frac{\mathfrak B^{(i)}}{hh+2h}+\tfrac12\mathfrak A^{(i)}\left(\frac1{h+2}-\frac{4h}{hh-1}+\frac3h\right)\right\}\\
&=\frac{aae}{h(h+2)}\left\{\begin{aligned}
&\tfrac12C^{(i)}-\tfrac12B^{(i)}-\frac{3A^{(i)}}{1-\theta}\\
&-2\mathfrak A^{(i)}\left(\frac{3+ih(h-1)}{hh-1}\right)
\end{aligned}\right\}\mathpunct{\cdot}
\end{aligned}
\]

VIII) In \(\delta r\) the coefficient of \(\cos((i+1)D+M')\), [which] \(aae'D^{(i)}\) \hypertarget{native-400}{}[may denote], is given by
\[
\begin{aligned}
aae'D^{(i)}&=\frac{aae'\mathfrak C^{(i)}}{hh+2h}\\
&=\frac{aae'}{hh+2h}\left\{\begin{aligned}
&\frac{i+1}{h+1}\bigl((2i+1)A^{(i+1)}-B^{(i+1)}\bigr)\\
&+(i+\tfrac12)B^{(i+1)}-\tfrac12C^{(i+1)}
\end{aligned}\right\}\mathpunct{\cdot}
\end{aligned}
\]

[Here too, the arguments \(iD+M\) and \((i+1)D+M'\) differ only by constant quantities. The terms can also be combined into a single term. If, proceeding as in VI, we set]
\[
\alpha\genfrac{}{}{0pt}{1}{\sin}{\cos}(x-A)+\beta\genfrac{}{}{0pt}{1}{\sin}{\cos}(x-B)
=f\genfrac{}{}{0pt}{1}{\sin}{\cos}(x-A+z),
\]
[then \(f\) and \(z\) are determined from
\[
\begin{aligned}
\beta\sin(A-B)&=f\sin z\\
\alpha+\beta\cos(A-B)&=f\cos z.]
\end{aligned}
\]

IX) In \(\delta v\) the coefficient of \(\sin(iD+M)\) becomes
\[
\begin{aligned}
&=ae\left\{\begin{aligned}
&-2(h+1)E^{(i)}+\frac{2\mathfrak B^{(i)}-\tfrac12\beta^{(i)}}{h+1}\\
&-(2h+1)\frac{\mathfrak A^{(i)}}{hh-1}
\end{aligned}\right\}\\
&=\frac{ae}{h+1}\left\{\begin{aligned}
&-2E^{(i)}-2\mathfrak A^{(i)}\frac{\tfrac12i(h-1)-3}{hh-1}\\
&+\frac{(i+1)A^{(i)}}{1-\theta}
\end{aligned}\right\}\mathpunct{\cdot}
\end{aligned}
\]

X) In \(\delta v\) the coefficient of \(\sin((i+1)D+M')\) becomes
\[
\begin{aligned}
&=ae'\left\{-2(h+1)D^{(i)}+\frac{2\mathfrak C^{(i)}-\tfrac12\gamma^{(i)}}{h+1}\right\}\\
&=\frac{ae'}{h+1}\left\{\begin{aligned}
&-2D^{(i)}\\
&-\tfrac12\frac{i+1}{h+1}\bigl((2i+1)A^{(i+1)}-B^{(i+1)}\bigr)
\end{aligned}\right\}\mathpunct{\cdot}
\end{aligned}
\]
\par\medskip\centerline{\rule{.23\textwidth}{.35pt}}\medskip
